\section*{Capítulo \ref{ch:g10:why-clouds-form} --- Por que as nuvens se formam}

\begin{solution}{exo:g10:why-clouds-form:1}
A \omterm{def:g10:why-clouds-form:humidity}{umidade relativa} é o teor de vapor de água do ar dividido pelo seu teor de
\omterm{def:g10:why-clouds-form:humidity}{saturação} à mesma temperatura e à mesma pressão, em porcentagem. O \omterm{def:g10:why-clouds-form:dew-point}{ponto de
orvalho} é a temperatura até a qual o ar precisa ser resfriado, conservando o
seu teor de água e a sua pressão, para ficar \omterm{def:g10:why-clouds-form:humidity}{saturado}.
\end{solution}

\begin{solution}{exo:g10:why-clouds-form:2}
\num{3.8}, \num{7.6}, \num{14.7} e \qty{27.2}{g/kg}. Cada
\qty{10}{\degreeCelsius} o multiplica por cerca de 2 ($7.6/3.8 = 2.0$,
$14.7/7.6 \approx 1.9$, $27.2/14.7 \approx 1.9$).
\end{solution}

\begin{solution}{exo:g10:why-clouds-form:3}
Teor de \omterm{def:g10:why-clouds-form:humidity}{saturação} a \qty{24}{\degreeCelsius}: \qty{18.9}{g/kg}; \omterm{def:g10:why-clouds-form:humidity}{umidade
relativa} $10.0/18.9 \approx \qty{53}{\percent}$. O teor de \omterm{def:g10:why-clouds-form:humidity}{saturação} é de
\qty{10.0}{g/kg} a \qty{14}{\degreeCelsius}: é o \omterm{def:g10:why-clouds-form:dew-point}{ponto de orvalho}.
\end{solution}

\begin{solution}{exo:g10:why-clouds-form:4}
No verão, o ar é úmido e o seu \omterm{def:g10:why-clouds-form:dew-point}{ponto de orvalho} fica bem acima da
temperatura do copo: o ar que toca o copo se resfria abaixo do seu \omterm{def:g10:why-clouds-form:dew-point}{ponto de
orvalho} e a água se condensa sobre ele. O ar de inverno muito seco dentro de
casa tem um \omterm{def:g10:why-clouds-form:dew-point}{ponto de orvalho} mais baixo que a temperatura do copo: mesmo
resfriado contra ele, o ar continua não \omterm{def:g10:why-clouds-form:humidity}{saturado}.
\end{solution}

\begin{solution}{exo:g10:why-clouds-form:5}
Aquecimento por baixo (cúmulos da tarde sobre um chão ensolarado); montanhas
(vento forçado a subir uma cordilheira); frentes (ar quente subindo sobre ar
frio); convergência (ventos que convergem perto do equador ou para dentro de
uma depressão).
\end{solution}

\begin{solution}{exo:g10:why-clouds-form:6}
$125 \times (30 - 18) = \qty{1500}{m}$; $125 \times (22 - 20) =
\qty{250}{m}$. O segundo dia: as nuvens se formam a apenas \qty{250}{m}
acima do chão.
\end{solution}

\begin{solution}{exo:g10:why-clouds-form:7}
A \qty{1}{km}: $18 - 9.8 = \qty{8.2}{\degreeCelsius}$; a \qty{2}{km}:
$18 - 19.6 = \qty{-1.6}{\degreeCelsius}$. Descendo, ela se aquece à mesma
taxa e chega ao chão de novo a \qty{18}{\degreeCelsius}.
\end{solution}

\begin{solution}{exo:g10:why-clouds-form:8}
Com \qty{4}{\degreeCelsius/km}, tanto uma parcela não \omterm{def:g10:why-clouds-form:humidity}{saturada}
(\qty{9.8}{\degreeCelsius/km}) quanto uma \omterm{def:g10:why-clouds-form:humidity}{saturada}
(\qty{5}{\degreeCelsius/km}) se resfriam mais depressa que o ar em volta:
\omterm{def:g10:why-clouds-form:stability}{estável} nos dois casos. Com \qty{12}{\degreeCelsius/km}, até uma bolha não
\omterm{def:g10:why-clouds-form:humidity}{saturada} se resfria mais devagar que o ar em volta: ela continua mais quente
e segue subindo, e as bolhas de ar quente sobem livremente.
\end{solution}

\begin{solution}{exo:g10:why-clouds-form:9}
$14.7 - 7.6 = \qty{7.1}{g}$ por quilograma de ar. O \omterm{def:g10:earth-energy-budget:latent-heat}{calor latente} aquece o
ar em que a água se condensa, o que retarda o seu resfriamento.
\end{solution}

\begin{solution}{exo:g10:why-clouds-form:10}
Os volumes crescem como o cubo do diâmetro: $(0.5/0.02)^3 = 25^3 = 15\,625$
gotículas para uma gota de garoa; $(2/0.02)^3 = 100^3 = 1\,000\,000$ para
uma gota de chuva.
\end{solution}

\begin{solution}{exo:g10:why-clouds-form:11}
Sob um céu limpo, o chão irradia infravermelho e se resfria, assim como o ar
que o toca. Esse ar frio é mais denso que o ar acima dele; nas encostas, ele
escorre morro abaixo e se acumula no fundo do vale. Ali ele continua se
resfriando durante a noite até atingir o seu \omterm{def:g10:why-clouds-form:dew-point}{ponto de orvalho}, e gotículas
se formam no próprio ar: a neblina. No alto dos morros, o ar não fica tão
frio e permanece acima do seu \omterm{def:g10:why-clouds-form:dew-point}{ponto de orvalho}.
\end{solution}

\begin{solution}{exo:g10:why-clouds-form:12}
Na altura $z$, a temperatura é $T - a z$ e o \omterm{def:g10:why-clouds-form:dew-point}{ponto de orvalho} $T_d - b z$;
eles são iguais quando $(a - b) z = T - T_d$, logo $z = (T - T_d)/(a - b)$.
Com $b = 1.7$: $a - b = 8.1$, e $1000/8.1 \approx \qty{123}{m}$ por grau.
\end{solution}

\begin{solution}{exo:g10:why-clouds-form:13}
\omterm{def:g10:why-clouds-form:condensation-level}{Nível de condensação}: $125 \times (15 - 11) = \qty{500}{m}$; temperatura
ali: $15 - 9.8 \times 0.5 \approx \qty{10.1}{\degreeCelsius}$. Crista: $10.1 - 5
\times 1.5 = \qty{2.6}{\degreeCelsius}$. Descendo \qty{2}{km} segundo o
gradiente seco: $2.6 + 9.8 \times 2 \approx \qty{22.2}{\degreeCelsius}$,
cerca de sete graus mais quente do que do lado de barlavento.
\end{solution}

\begin{solution}{exo:g10:why-clouds-form:14}
A \qty{-20}{\degreeCelsius}: $126/103 \approx 1.22$, \qty{22}{\percent} de
supersaturação para o gelo; a \qty{-10}{\degreeCelsius}: $287/260 \approx
1.10$, \qty{10}{\percent}. Nessas temperaturas, quanto mais fria a nuvem,
maior a vantagem do gelo: o processo funciona melhor na parte alta e fria
de uma nuvem onde gotículas e cristais estão misturados.
\end{solution}

\begin{solution}{exo:g10:why-clouds-form:15}
Não \omterm{def:g10:why-clouds-form:humidity}{saturada}: a parcela se resfria \qty{9.8}{\degreeCelsius}, o ar em volta
\qty{6.5}{\degreeCelsius}: ela fica \qty{3.3}{\degreeCelsius} mais fria e
volta a descer. \omterm{def:g10:why-clouds-form:humidity}{Saturada}: ela se resfria \qty{5}{\degreeCelsius}, logo fica
\qty{1.5}{\degreeCelsius} mais quente e continua subindo. O ar só é \omterm{def:g10:why-clouds-form:stability}{instável}
sob uma condição: que o ar que sobe esteja \omterm{def:g10:why-clouds-form:humidity}{saturado}.
\end{solution}

\begin{solution}{pb:g10:why-clouds-form:1}
\textbf{1.} \qty{14.7}{g/kg}.

\textbf{2.} $11.4/14.7 \approx \qty{78}{\percent}$.

\textbf{3.} \qty{16}{\degreeCelsius}, em que o teor de \omterm{def:g10:why-clouds-form:humidity}{saturação} é de
\qty{11.4}{g/kg}.

\textbf{4.} $11.4/24.1 \approx \qty{47}{\percent}$.

\textbf{5.} Não: ele contém os mesmos \qty{11.4}{g/kg}. A sua \omterm{def:g10:why-clouds-form:humidity}{umidade
relativa} caiu porque, mais quente, ele poderia conter mais.

\textbf{6.} No seu \omterm{def:g10:why-clouds-form:dew-point}{ponto de orvalho}, \qty{16}{\degreeCelsius}.

\textbf{7.} $11.4 - 8.7 = \qty{2.7}{g}$ por quilograma de ar.

\textbf{8.} Sob um céu noturno limpo, a grama e os tetos irradiam
infravermelho e ficam mais frios que o ar; o ar que os toca se resfria
abaixo do seu \omterm{def:g10:why-clouds-form:dew-point}{ponto de orvalho} antes do resto do ar.

\textbf{9.} $(28.0 - 13.4)/1.5 \approx \qty{9.7}{\degreeCelsius/km}$,
praticamente o gradiente seco: as bolhas da tarde misturaram essa camada.

\textbf{10.} $(13.4 - 0.4)/2.0 = \qty{6.5}{\degreeCelsius/km}$.

\textbf{11.} \omterm{def:g10:why-clouds-form:stability}{Estável} para uma parcela não \omterm{def:g10:why-clouds-form:humidity}{saturada}, que se resfria mais
depressa (\qty{9.8}{\degreeCelsius/km}); \omterm{def:g10:why-clouds-form:stability}{instável} para uma \omterm{def:g10:why-clouds-form:humidity}{saturada}, que se
resfria mais devagar (\qty{5}{\degreeCelsius/km}).

\textbf{12.} A temperatura aumenta com a altura, de \num{-0.3} para
\qty{1.2}{\degreeCelsius}: uma camada muito \omterm{def:g10:why-clouds-form:stability}{estável}, que funciona como uma
tampa. A nuvem que cresce para cima para ali: o seu topo fica perto de
\qty{4}{km}.

\textbf{13.} $28 - 16 = \qty{12}{\degreeCelsius}$.

\textbf{14.} Temperatura $28 - 9.8 = \qty{18.2}{\degreeCelsius}$; \omterm{def:g10:why-clouds-form:dew-point}{ponto de
orvalho} $16 - 1.8 = \qty{14.2}{\degreeCelsius}$.

\textbf{15.} De 12 para $18.2 - 14.2 = 4.0$: ela diminuiu
\qty{8}{\degreeCelsius}, uma taxa de \qty{8}{\degreeCelsius} por quilômetro.

\textbf{16.} Um grau de diferença se fecha em $1/8$ km, isto é,
\qty{125}{m}.

\textbf{17.} $125 \times (20 - 16) = \qty{500}{m}$.

\textbf{18.} O \omterm{def:g10:why-clouds-form:dew-point}{ponto de orvalho} fica em \qty{16}{\degreeCelsius} enquanto o
ar perto do chão se aquece: a diferença aumenta, e com ela a altura em que
ela se fecha.

\textbf{19.} $28 - 9.8\,z = 16 - 1.8\,z$.

\textbf{20.} $8\,z = 12$, logo $z = 1.5$: a base dos cúmulos da tarde fica a
\textbf{cerca de \qty{1.5}{km}} acima da planície (o cálculo exato com as
propriedades da água dá \qty{1.51}{km}).
\end{solution}
